\documentclass[11pt]{article}
\usepackage[a4paper,margin=25mm]{geometry}
\usepackage{amsmath,amssymb}
\usepackage[hidelinks]{hyperref}
\setlength{\emergencystretch}{3em}
\title{Conjecture 00000005045: local Hanf data do not give global isomorphism}
\author{lizaixi01, with OpenAI Codex assistance}
\date{23 October 2026}
\begin{document}
\maketitle
\section{The necessary transfer fails at degree two}
The source explicitly defines the local-to-global conversion as
conversion from Hanf local isomorphism to global isomorphism, and
asserts a uniform radius-two conversion constant table governed by
degree thresholds. Under that stated definition, the necessary
transfer already fails for finite graphs of degree two. For each
integer $t\geq0$, put
\[
 G_t=(t+1)C_{12},\qquad H_t=2(t+1)C_6,
\]
where multiplication denotes disjoint copies of the actual cycle
graph. Both graphs have $12(t+1)$ vertices, and every vertex has
degree $2$. Every rooted induced closed radius-two ball in either
graph is the same five-vertex path, rooted at its middle vertex.
Consequently the graphs have a bijection of vertices preserving
rooted radius-two neighborhood types. In fact, the number of roots
of every neighborhood type is exactly equal, so every finite Hanf
count threshold agrees. Yet $G_t$ and $H_t$ are not isomorphic.
Taking $t$ arbitrarily large excludes any size cutoff as well.

We do not assign numbers to the source's unspecified constant
table. Rather, any uniform radius-two threshold/size-cutoff rule
that accomplishes the stated conversion must work on this family,
and cannot do so. A table for a different, unstated operation is
outside this conclusion. No independent claim about numerical
entries $3$ or $5$, monotonicity, or a size-doubling jump is made.

\section{Actual balls, isomorphisms and Hanf counts}
Lean uses Mathlib's actual \texttt{SimpleGraph.cycleGraph}.
Disjoint copies have vertices $(i,v)$, with adjacency precisely
equality of copy indices and actual cycle adjacency. The closed
radius-two ball is defined by
\[
 B_2(v)=\{w:w=v\ \lor\ v\sim w\ \lor\
                 (\exists z,\ v\sim z\ \land\ z\sim w)\}.
\]
The proof formally identifies this with existence of an actual
graph walk of length at most $2$. It does not use the default
value of a distance between disconnected vertices. A rooted ball
isomorphism is a genuine graph isomorphism of the induced
subgraphs, carrying the actual root to the other root.

At root $0$, the vertices of these balls are respectively
$[0,1,2,10,11]$ in $C_{12}$ and $[0,1,2,4,5]$ in $C_6$.
The proof checks their actual adjacency relations, constructs the
induced graph isomorphism, and then uses genuine cycle rotations
and copy projections for arbitrary roots. Ordinary kernel-checked
finite decisions are used for the small base graphs.

\newpage
For any finite ambient graph $A$ and any rooted target graph
$(K,z)$, the number of roots whose induced ball is isomorphic
to $B_2^K(z)$ is the cardinality of that actual subtype of vertices.
A Hanf vertex bijection gives an explicit equivalence between
these subtypes in $G_t$ and $H_t$. Thus their finite cardinalities
are equal for every target, not merely for an assigned list of
counts. In particular, for every threshold $q$, the counts of
each finite rooted type are equal or both exceed $q$, with equality
always supplying the first alternative. \texttt{Nat.card} is used
only on these finite ambient vertex subtypes; no infinite-cardinality
default is invoked.

\section{Global nonisomorphism and arbitrary size}
Every actual walk in a disjoint-copy graph preserves the copy
index. The actual twelve-cycle is connected. If a global graph
isomorphism existed, the images of all twelve vertices of one
$C_{12}$ component would therefore lie in one $C_6$ component.
Injectivity of the graph isomorphism would induce an injection
$\operatorname{Fin}12\hookrightarrow\operatorname{Fin}6$, contradicting
$12>6$. This proves nonisomorphism for every $t$.

Lean proves the actual orders and all actual vertex degrees.
Its final theorem states that for every size cutoff $N$ and
every count threshold $q$, some $t$ gives equal-order graphs of
order at least $N$, degree exactly $2$, genuine Hanf matching
and threshold agreement, but no global isomorphism. The formal
negation of a uniform cutoff/threshold conversion follows.

\section{Scope, references and verification}
Hanf locality in its usual form yields bounded-quantifier-rank
first-order equivalence, not global graph isomorphism. This
submission does not refute that theorem. The distinction is
explicit in Anuj Dawar's primary Simons Institute lectures:
\emph{Finite and Algorithmic Model Theory II: Automata-Based
Methods}, 30 August 2016, PDF page 4, gives the neighborhood
bijection definition; lecture III, \emph{Parameterized
Satisfaction}, 31 August 2016, PDF pages 9--10, gives induced
neighborhoods, threshold counts and the first-order conclusion.
The directly read sources are
\href{https://simons.berkeley.edu/sites/default/files/docs/5233/lecture22.pdf}{lecture II}
and
\href{https://simons.berkeley.edu/sites/default/files/docs/5234/lecture31.pdf}{lecture III}.

The package pins Lean 4.33.0 and Mathlib commit
\begin{center}\small\texttt{db584cd6d46c92f209a44c0f1c829460d327499d}\end{center}
and supplies the original bilingual source, complete LaTeX and
PDF, all named local declaration axiom audits, a new root build
and full official \texttt{leanchecker --fresh Main} replay.
No \texttt{sorry}, \texttt{native\_decide} or additional axiom is used.
OpenAI Codex assisted, with separate supervising AI proof/source/
report/PDF review. This is not human certification; organizer
acceptance remains pending.
\end{document}
